\label{ch:inference-process}

This chapter documents the statistical inference machinery: the
gradient-based fit and its Laplace uncertainty, shared unchanged with the
rest of \texttt{tpeanuts} (\cref{sec:fit-procedure}); the correlated
Gaussian likelihood SNO's own analysis uses in place of a diagonal one
(\cref{sec:correlated-likelihood}); how the two covariance components are
built from the real published tables (\cref{sec:covariance-construction});
a genuine extension to this project's own shared fitting code added
specifically to support this document (\cref{sec:callable-extension}); and
the fixed-slice/profiled distinction for the 2-D likelihood scan
(\cref{sec:scan-modes}).

\section{Gradient-based fit procedure}
\label{sec:fit-procedure}

\modulehead{inference/fit.py}{\pyfunc{fit\_lbfgs}/\pyfunc{minimize\_lbfgs}:
LBFGS point estimate plus Laplace/Fisher covariance, shared unchanged by
every model in this project.}

\begin{theorybox}
$\hat\theta=\arg\min_\theta\bigl[-2\ln L(\theta)\bigr]$ is found by
\code{torch.optim.LBFGS} using the exact autograd gradient of the
statistic through every real physics step of
\cref{ch:detector-implementation} (no finite differences), after a
per-parameter rescaling that keeps $\theta_{12}\sim\mathcal O(1)$~rad and
$\Delta m^2_{21}\sim\mathcal O(10^{-5})$~eV$^2$ on comparable numerical
footing for LBFGS's own identity-initialised inverse-Hessian
approximation. The marginal uncertainty on each fitted parameter is the
Laplace approximation: with $\mathcal I=\tfrac12\nabla^2_\theta(-2\ln L)$
the Fisher information at the optimum (the autograd Hessian),
\begin{equation}
  \keyresult{\Sigma = \mathcal I^{-1}, \qquad \sigma_j = \sqrt{\Sigma_{jj}}.}
  \label{eq:laplace-covariance}
\end{equation}
A direction the data barely constrains gives a Hessian eigenvalue near
zero (occasionally slightly negative from numerical round-off); small or
negative eigenvalues are floored at a tiny fraction of the largest one
before inversion, so an unconstrained direction gets a large but finite
$\sigma_j$ instead of \code{nan} -- exactly the behaviour
\cref{sec:results-fit-total} documents for $\Delta m^2_{21}$ once
$V_{\rm syst}$ is included.
\end{theorybox}

\begin{implbox}
\pyfunc{fit\_lbfgs(model, theta0, value, sigma\_minus=None,
sigma\_plus=None, likelihood=..., max\_iter=200)} returns a
\pyclass{FitResult} (\code{theta\_hat}, \code{param\_names},
\code{chi2\_history}, \code{covariance}, and a \code{sigma} property
implementing \labelcref{eq:laplace-covariance}). Identical machinery to
every other detector in this project (Daya Bay, Borexino, SNO Phase~I,
IceCube); this document's only change to it is
\cref{sec:callable-extension}'s \code{str | Callable} dispatch extension.
\end{implbox}

\section{Correlated Gaussian likelihood}
\label{sec:correlated-likelihood}

\modulehead{inference/likelihood.py}{\pyfunc{correlated\_gaussian\_nll},
\pyfunc{cholesky\_from\_covariance}.}

\begin{theorybox}
SNO's own oscillation analysis (Section~XIII, Eq.~19-21 of
\cite{Aharmim2005SaltPhase}) compares the real data and model 38-vectors
with a full \emph{correlated} $\chi^2$, not the diagonal one every other
model in this project uses:
\begin{equation}
  \keyresult{\chi^2 = \sum_{i,j=1}^{38}\bigl(Y_i^{\rm data}-Y_i^{\rm model}\bigr)
    \bigl[\sigma_{ij}^2({\rm tot})\bigr]^{-1}\bigl(Y_j^{\rm data}-Y_j^{\rm model}\bigr),
    \qquad \sigma^2({\rm tot}) = \sigma^2({\rm stat}) + \sigma^2({\rm syst}).}
  \label{eq:correlated-chi2}
\end{equation}
Rather than explicitly inverting the $38\times38$ covariance -- numerically
the less stable of the two standard constructions -- the Cholesky factor
$L$ ($\sigma^2({\rm tot})=LL^T$) is used instead,
\begin{equation}
  \chi^2 = \bigl\|L^{-1}(\vec Y^{\rm data}-\vec Y^{\rm model})\bigr\|^2,
  \label{eq:cholesky-chi2}
\end{equation}
solved as a triangular linear system rather than a matrix inversion. With a
diagonal covariance, \labelcref{eq:cholesky-chi2} reduces exactly to this
project's existing \pyfunc{chi2\_asymmetric} statistic with
$\sigma_{\rm minus}=\sigma_{\rm plus}=\sqrt{{\rm diag}(\sigma^2)}$ -- the
two-sided special case of the one-sided general form -- verified directly
in \texttt{inference/test/test1\_likelihood.py}.
\end{theorybox}

\begin{implbox}
\pyfunc{correlated\_gaussian\_nll(prediction, value, sigma\_minus=None,
sigma\_plus=None, *, cholesky\_L)} implements
\labelcref{eq:cholesky-chi2} via \texttt{torch.linalg.solve\_triangular};
\code{sigma\_minus}/\code{sigma\_plus} are accepted but must be
\code{None}, purely so this function shares
\pyfunc{chi2\_asymmetric}/\pyfunc{poisson\_nll}'s call signature (needed
for the dispatch extension, \cref{sec:callable-extension}).
\pyfunc{cholesky\_from\_covariance(covariance)} is a thin
\texttt{torch.linalg.cholesky} wrapper, meant to be called \emph{once},
outside any fit or scan loop, whenever the covariance itself does not
depend on the free parameters -- exactly the case here, since
$\sigma^2({\rm tot})$ is a fixed, published quantity.
\end{implbox}

\begin{resultbox}
Verified on a random dense $5\times5$ positive-definite matrix: the
Cholesky-solved $\chi^2$ agrees with the textbook explicit-inverse form
$\vec r^T V^{-1}\vec r$ to within $10^{-8}$ absolute -- confirming
\labelcref{eq:cholesky-chi2} is the same computation as
\labelcref{eq:correlated-chi2}, not an approximation to it.
\end{resultbox}

\section{Building the covariance from the real published tables}
\label{sec:covariance-construction}

\modulehead{detector/sno\_ii/io.py}{\pyfunc{load\_observed\_vector\_and\_covariance},
\pyfunc{build\_systematic\_covariance},
\pyfunc{load\_observed\_vector\_and\_total\_covariance}.}

\begin{theorybox}
Day and night are statistically independent in SNO's own signal extraction
(they come from separate fits to separate data subsets), so the
statistical component of \labelcref{eq:correlated-chi2} is exactly
block-diagonal,
\begin{equation}
  \keyresult{\sigma^2({\rm stat}) = \begin{pmatrix} V_{\rm stat}^{\rm day} & 0 \\ 0 & V_{\rm stat}^{\rm night}\end{pmatrix},
    \qquad V_{\rm stat}^{p} = {\rm diag}(\vec\sigma^p)\,C^p\,{\rm diag}(\vec\sigma^p),}
  \label{eq:stat-covariance}
\end{equation}
with $C^p$ the real published $19\times19$ correlation matrix for period
$p$ (Tables XXXII/XXXIII) and $\vec\sigma^p$ the real per-channel
statistical uncertainties (Tables XXX/XXIV).

The systematic component (Eq.~21 of \cite{Aharmim2005SaltPhase}) is built
from Table XXXIV's per-bin percent sensitivities to 16 detector
systematics,
\begin{equation}
  \keyresult{\sigma_{ij}^2({\rm syst}) = \sum_{k=1}^{16}\frac{\partial Y_i}{\partial S_k}
    \frac{\partial Y_j}{\partial S_k}(\Delta S_k)^2,
    \qquad \frac{\partial Y_i}{\partial S_k} = \frac{{\rm pct}_{ik}}{100}\,Y_i,}
  \label{eq:syst-covariance}
\end{equation}
with $\Delta S_k=1$ (the tabulated percentages already are per-$1\sigma$
sensitivities) and $Y_i$ each observable's own real central value. Unlike
$\sigma^2({\rm stat})$, this component is \emph{not} block-diagonal in
day/night: the same physical systematic (e.g.\ energy scale) affects both
periods' data identically, since it is the same instrument. Table XXXIV's
own plus/minus asymmetric percentages are symmetrized (averaged in
absolute value) before use -- Eq.~\labelcref{eq:syst-covariance}'s simple
quadratic form has no asymmetric-uncertainty counterpart.
\end{theorybox}

\begin{implbox}
\pyfunc{load\_observed\_vector\_and\_covariance} implements
\labelcref{eq:stat-covariance}. \pyfunc{build\_systematic\_covariance(value)}
implements \labelcref{eq:syst-covariance} vectorized as a single Gram
product, $\sigma^2({\rm syst}) = \Delta Y^T\Delta Y$ with $\Delta Y$ shaped
$(16,38)$ -- symmetric and positive-semi-definite by construction, with no
explicit symmetry check required. \pyfunc{load\_observed\_vector\_and\_total\_covariance}
returns $(\vec Y^{\rm data}, \sigma^2({\rm stat})+\sigma^2({\rm syst}))$
directly.
\end{implbox}

\begin{resultbox}
Both $19\times19$ statistical matrices and the full $38\times38$ combined
covariance are verified symmetric and (strictly, not just semi-)
positive-definite before every fit in \cref{ch:results}, with minimum
eigenvalues $6.25\times10^{-6}$ ($V_{\rm stat}$ alone) and
$6.37\times10^{-6}$ ($V_{\rm stat}+V_{\rm syst}$) -- comfortably away from
zero, so \texttt{torch.linalg.cholesky} never fails.
\end{resultbox}

\section{Extending the shared fit machinery: \texorpdfstring{\code{str | Callable}}{str or Callable} likelihoods}
\label{sec:callable-extension}

\modulehead{inference/fit.py, inference/scan.py}{\pyfunc{\_resolve\_likelihood}.}

\begin{theorybox}
Every likelihood used elsewhere in this project (\pyfunc{chi2\_asymmetric},
\pyfunc{poisson\_nll}) is looked up in a fixed string-keyed registry,
\code{LIKELIHOODS}, inside \pyfunc{fit\_lbfgs}/\pyfunc{minimize\_lbfgs}/
\pyfunc{loglik\_grid}. \pyfunc{correlated\_gaussian\_nll} cannot be
registered this way: it needs a mandatory \code{cholesky\_L} argument with
no sensible default, so a bare string lookup would immediately fail with a
missing-argument error. Rather than special-case this one likelihood
inside every fitting function, this document's own work extended the
\code{likelihood} argument of \pyfunc{fit\_lbfgs},
\pyfunc{minimize\_lbfgs}, and \pyfunc{loglik\_grid} to accept
\emph{either} a registry string \emph{or} a directly-supplied callable
sharing the same \code{(prediction, value, sigma\_minus, sigma\_plus) ->
loss} call shape -- the covariance-specific argument is bound in advance
via \code{functools.partial(correlated\_gaussian\_nll, cholesky\_L=L)},
and the bound partial is passed exactly where a string would go.
\end{theorybox}

\begin{implbox}
\pyfunc{\_resolve\_likelihood(likelihood)}: if \code{likelihood} is a
string, looks it up in \code{LIKELIHOODS} (raising the same
\code{ValueError} as before on an unknown name); otherwise returns it
unchanged. Every call site that previously did
\code{loss\_fn = LIKELIHOODS[likelihood]} now calls this helper instead --
a change confined to \emph{how} the loss function is resolved, not to any
of the surrounding fit/scan logic, so every existing detector's own
string-based calls (Daya Bay, Borexino, SNO Phase~I, IceCube) are
unaffected. \pyfunc{calibrate\_delta\_threshold} was deliberately left
requiring a string only: its own Monte-Carlo toy-sampling logic
(\code{\_sample\_toy\_value}) does not yet know how to draw a correlated
multivariate-Gaussian toy, so passing a callable there fails immediately
and legibly rather than silently sampling incorrectly.
\end{implbox}

\begin{resultbox}
Verified end to end on a synthetic linear model
(\texttt{inference/test/test1\_likelihood.py}): \pyfunc{fit\_lbfgs} with a
bound \pyfunc{correlated\_gaussian\_nll} callable recovers a known,
noiseless linear truth to within $10^{-6}$ absolute; \pyfunc{loglik\_grid}
with the same callable locates the correct grid minimum. Both behave
identically to the pre-existing string-based dispatch path.
\end{resultbox}

\section{Fixed-slice vs.\ profiled likelihood scans}
\label{sec:scan-modes}

\modulehead{inference/scan.py}{\pyfunc{loglik\_grid}.}

\begin{theorybox}
On a chosen $(\theta_{12},\Delta m^2_{21})$ grid, $\lambda_{8B}$ can either
be held fixed at its own best-fit value from a prior full fit (a cheap
slice through the likelihood surface, one forward pass per grid point) or
re-minimized independently at every grid point (\code{profile\_others=True},
the statistically correct construction whenever the profiled parameter is
itself uncertain, at the cost of one LBFGS run per grid point). The two
constructions are not generally interchangeable: profiling out an
uncertain nuisance parameter can only widen a resulting confidence region
relative to holding it fixed, never shrink it, since the fixed-slice
surface is by definition an upper bound on the profiled one at every grid
point (the profiled value is chosen to minimize the statistic, the fixed
value generally does not).
\end{theorybox}

\begin{implbox}
\pyfunc{loglik\_grid(model, theta\_baseline, param\_x, x\_grid, param\_y,
y\_grid, value, likelihood=..., profile\_others=False)} implements both
modes via its own \pyclass{\_FixedSliceModel} adapter, which re-exposes
every free parameter other than the two scanned ones to
\pyfunc{minimize\_lbfgs} at each grid point when
\code{profile\_others=True}. No new scan code was written for this
document -- \cref{ch:results}'s profiled scan reuses this pre-existing
machinery directly, unmodified beyond \cref{sec:callable-extension}'s
\code{str | Callable} extension.
\end{implbox}
